\section{Story Explorer - AI Evaluierung}
Siehe \href{sec:KI}{KI Nutzung} für details, Link zum \href{https://gemini.google.com/share/354e7d9f76b8}{Prompt}\\

\noindent
Hallo! Als jemand, der sich für städtische Themen in Wien interessiert, aber keinen technischen Hintergrund in der Datenanalyse hat, finde ich die Grundidee eurer „Story Explorer“-App wirklich spannend.
\\

\noindent
Was mir besonders gefällt:
\\

\noindent
    „Kurz erklärt“-Ansatz: Dass nackte Daten nicht einfach nur hingeworfen, sondern in einfacher Sprache eingeordnet werden (wie beim PM2.5-Beispiel), holt Laien wie mich perfekt ab.
\\

\noindent
    Barrierefreiheit: Dass dieser Aspekt (Farbschemata, Darstellungsgrößen) so prominent bedacht wurde, ist bei Angeboten rund um öffentliche Daten essenziell.
\\

\noindent
    Kognitive Entlastung: Der lineare Flow (Thema → Datensatz → Story) und die Beschränkung auf vier Kernthemen wirken sehr aufgeräumt. Man wird nicht von der schieren Menge des Open-Data-Katalogs erschlagen.
\\

\noindent
Wo ich noch Verbesserungspotenzial sehe:
\\

\noindent
    Usability in der Übersicht (Screen 2): Die Auswahl der Datensätze über Radiobuttons wirkt für eine moderne, mobile Touch-App ungewohnt und etwas unpraktisch. Intuitiver wäre es, wenn einfach die gesamte „Karte“ (Card) des jeweiligen Eintrags klickbar wäre, um zur Ansicht zu gelangen.
\\

\noindent
    Begrifflichkeit „Story“: Der Name verspricht Datenstorys. Der Prototyp in Screen 3 zeigt jedoch aktuell „nur“ ein klassisches Balkendiagramm mit Erklärtext. Bei einer Story erwarte ich fast schon ein narratives Element, vielleicht Erkenntnisse, durch die man swipen kann. Hier solltet ihr überlegen, ob die Visualisierung interaktiver wird oder ob der Begriff „Story“ vielleicht eine falsche Erwartung weckt.
\\

\noindent
    Personalisierung: In der Navigationsleiste gibt es ein Lesezeichen-Symbol („Saved“), das im Text aber nicht näher beschrieben wird. Als Nutzer fände ich es zudem super, wenn ich in den Einstellungen meinen Wohnbezirk fest hinterlegen könnte, sodass dieser bei allen Abfragen als Standardfilter dient.
\\

\noindent
Fazit:
Ein tolles Konzept, um die Wiener Open Data in den Alltag der Bürger:innen zu bringen. Wenn die Bedienung noch etwas touch-freundlicher wird, würde ich die App definitiv nutzen!